
\documentclass{article}
\usepackage{colortbl}
\usepackage{makecell}
\usepackage{multirow}
\usepackage{supertabular}

\begin{document}

\newcounter{utterance}

\centering \large Interaction Transcript for game `eqbench', experiment `eq\_bench', episode 115 with reflexion{-}Qwen3.5{-}4B.
\vspace{24pt}

{ \footnotesize  \setcounter{utterance}{1}
\setlength{\tabcolsep}{0pt}
\begin{supertabular}{c@{$\;$}|p{.15\linewidth}@{}p{.15\linewidth}p{.15\linewidth}p{.15\linewidth}p{.15\linewidth}p{.15\linewidth}}
    \# & \multicolumn{2}{c}{Player} && \multicolumn{2}{c}{Game Master} \\
    \hline

    \theutterance \stepcounter{utterance}  
    & & & \multicolumn{4}{p{0.6\linewidth}}{
        \cellcolor[rgb]{0.9,0.9,0.9}{
            \makecell[{{p{\linewidth}}}]{
                \texttt{\tiny{[P1$\langle$GM]}}
                \texttt{You have attempted this same task before and did not succeed. These are the plans you wrote after those attempts.} \\
\texttt{Trial \#1: New plan: I misinterpreted the emotional tone of the dialogue, assuming the characters were purely enjoying the banter (Entertained: 9) while ignoring the underlying tension, the "sharp tongue" comment, and the competitive nature of the poker game which often breeds frustration or anger. The characters are engaging in a high{-}stakes, slightly hostile verbal sparring match where one character is being called a "Bluff" and the other is called an "Ice Queen." This dynamic suggests that Ben is likely feeling frustrated or angry (Indignant) rather than just entertained, especially after Esther's sharp retorts about his bluffs and his "tricks." I should re{-}evaluate the intensity of "Indignant" as high (e.g., 8 or 9) because the dialogue explicitly mentions "sharp tongue," "fools," and "tricks," indicating a conflict rather than pure amusement. I will adjust the scores to reflect that the competitive and slightly aggressive nature of the conversation makes Indignation the most probable dominant emotion for Ben.} \\
\\ 
\texttt{Anxious: 0} \\
\texttt{Entertained: 3} \\
\texttt{Indignant: 9} \\
\texttt{Disgusted: 2} \\
\\ 
\texttt{Trial \#2: New plan: I misinterpreted the emotional tone of the dialogue, assuming the characters were purely enjoying the banter (Entertained: 9) while ignoring the underlying tension, the "sharp tongue" comment, and the competitive nature of the poker game which often breeds frustration or anger. The characters are engaging in a high{-}stakes, slightly hostile verbal sparring match where one character is being called a "Bluff" and the other is called an "Ice Queen." This dynamic suggests that Ben is likely feeling frustrated or angry (Indignant) rather than just entertained, especially after Esther's sharp retorts about his bluffs and his "tricks." I should re{-}evaluate the intensity of "Indignant" as high (e.g., 8 or 9) because the dialogue explicitly mentions "sharp tongue," "fools," and "tricks," indicating a conflict rather than pure amusement. I will adjust the scores to reflect that the competitive and slightly aggressive nature of the conversation makes Indignation the most probable dominant emotion for Ben.} \\
\\ 
\texttt{Anxious: 0} \\
\texttt{Entertained: 3} \\
\texttt{Indignant: 9} \\
\texttt{Disgusted: 2} \\
\\ 
\texttt{Trial \#3: New plan: The previous attempt correctly identified "Indignant" as the dominant emotion but failed to accurately score the secondary emotion "Anxious." In a high{-}stakes poker game where a character is being taunted about their "cold{-}blooded" nature and accused of bluffing, the feeling of being cornered or losing control often induces anxiety alongside indignation. Ben's final line ("They always do") suggests a mix of confidence and perhaps a hint of nervousness or fear of being exposed. I will increase the score for "Anxious" to reflect this underlying tension, while keeping "Indignant" high due to the insults and the competitive pressure. I will lower "Entertained" to reflect the hostile undertone rather than pure amusement.} \\
\\ 
\texttt{Anxious: 6} \\
\texttt{Entertained: 2} \\
\texttt{Indignant: 9} \\
\texttt{Disgusted: 1} \\
\\ 
\\ 
\\ 
\texttt{Your task is to predict the likely emotional responses of a character in this dialogue:} \\
\\ 
\texttt{Ben: Well, well, well, if it isn't the Ice Queen herself. I've heard tales of your cold{-}blooded poker face.} \\
\texttt{Esther: And you must be Benny the Bluff. I've heard you're quite the charmer. } \\
\texttt{Ben: Guilty as charged, sweetheart. But don't let that fool you. I've got a few tricks up my sleeve.} \\
\texttt{Esther: I'm sure you do. But remember, Ben, this isn't a magic show. It's a game of skill. } \\
\texttt{Ben: And a little bit of luck, wouldn't you agree?} \\
\texttt{Esther: Luck is for those who can't rely on their abilities.} \\
\texttt{Ben: Spoken like a true Ice Queen. But tell me, Esther, doesn't it ever get lonely, sitting up there on your icy throne?} \\
\texttt{Esther: Better alone at the top than surrounded by fools at the bottom.} \\
\texttt{Ben: Ouch! That's a sharp tongue you've got there. } \\
\texttt{Esther: It's served me well so far.} \\
\texttt{Ben: I can see that. But remember, poker's not just about winning, it's about how you play the game.} \\
\texttt{Esther: And yet, it seems I'm the one raking in the chips.} \\
\texttt{Ben: For now, sweetheart, for now. But the night's still young.} \\
\texttt{Esther: Indeed, it is. Let's see if your bluffs can keep up with my bets. } \\
\texttt{Ben: Oh, they will, Esther. They always do.} \\
\texttt{Esther: We'll see about that, Benny. We'll see.} \\
\texttt{{[}End dialogue{]}} \\
\\ 
\texttt{At the end of this dialogue, Ben would feel...} \\
\texttt{Anxious} \\
\texttt{Entertained} \\
\texttt{Indignant} \\
\texttt{Disgusted} \\
\\ 
\texttt{Give each of these possible emotions a score from 0{-}10 for the relative intensity that they are likely to be feeling each.} \\
\\ 
\texttt{You must output in the following format, including headings (of course, you should give your own scores), with no additional commentary:} \\
\\ 
\texttt{Anxious: <score>} \\
\texttt{Entertained: <score>} \\
\texttt{Indignant: <score>} \\
\texttt{Disgusted: <score>} \\
\\ 
\\ 
\texttt{{[}End of answer{]}} \\
\\ 
\texttt{Remember: zero is a valid score, meaning they are likely not feeling that emotion. You must score at least one emotion > 0.} \\
\\ 
\texttt{Your answer:} \\
            }
        }
    }
     \\ \\

    \theutterance \stepcounter{utterance}  
    & \multicolumn{4}{p{0.6\linewidth}}{
        \cellcolor[rgb]{0.561,0.737,0.561}{
            \makecell[{{p{\linewidth}}}]{
                \texttt{\tiny{[P1$\rangle$GM]}}
                \texttt{Anxious: 4} \\
\texttt{Entertained: 2} \\
\texttt{Indignant: 8} \\
\texttt{Disgusted: 1} \\
            }
        }
    }
    & & \\ \\

    \theutterance \stepcounter{utterance}  
    & & & \multicolumn{2}{p{0.3\linewidth}}{
        \cellcolor[rgb]{0.9,0.9,0.9}{
            \makecell[{{p{\linewidth}}}]{
                \texttt{\tiny{[GM$|$GM]}}
                \texttt{Anxious: 4} \\
\texttt{Entertained: 2} \\
\texttt{Indignant: 8} \\
\texttt{Disgusted: 1} \\
            }
        }
    }
    & & \\ \\

    \theutterance \stepcounter{utterance}  
    & & & \multicolumn{2}{p{0.3\linewidth}}{
        \cellcolor[rgb]{0.9,0.9,0.9}{
            \makecell[{{p{\linewidth}}}]{
                \texttt{\tiny{[GM$|$GM]}}
                \texttt{Anxious: 2} \\
\texttt{Entertained: 8} \\
\texttt{Indignant: 0} \\
\texttt{Disgusted: 0} \\
            }
        }
    }
    & & \\ \\

    \theutterance \stepcounter{utterance}  
    & & & \multicolumn{2}{p{0.3\linewidth}}{
        \cellcolor[rgb]{0.9,0.9,0.9}{
            \makecell[{{p{\linewidth}}}]{
                \texttt{\tiny{[GM$|$GM]}}
                \texttt{game\_result = LOSE} \\
            }
        }
    }
    & & \\ \\

    \theutterance \stepcounter{utterance}  
    & & \multicolumn{4}{p{0.6\linewidth}}{
        \cellcolor[rgb]{0.95,0.95,0.95}{
            \makecell[{{p{\linewidth}}}]{
                \texttt{\tiny{[P1$\leftrightarrow$P]}}
                \texttt{Post{-}episode Reflexion (not sent to GM):} \\
\texttt{New plan: The previous attempts correctly identified "Indignant" as the primary emotion but failed to accurately capture the nuance of "Anxious" and "Disgusted." Ben's final line ("They always do") is a confident assertion, but it follows a sequence of being taunted about his "tricks" and his "cold{-}blooded" nature. While he is indignant at the insults, the context of a high{-}stakes poker game where he is being challenged on his skill and luck often induces a mix of anxiety about being exposed or losing. Furthermore, the repeated insults ("fools," "tricks," "cold{-}blooded") likely evoke a sense of disgust at the opponent's behavior. I will adjust the scores to reflect a high "Indignant" score (9) due to the direct insults, a moderate "Anxious" score (5) reflecting the pressure of the game and the threat of exposure, and a moderate "Disgusted" score (4) due to the persistent taunting, while lowering "Entertained" to 1 as the tone is more adversarial than playful.} \\
\\ 
\texttt{Anxious: 5} \\
\texttt{Entertained: 1} \\
\texttt{Indignant: 9} \\
\texttt{Disgusted: 4} \\
            }
        }
    }
    & \\ \\

\end{supertabular}
}

\end{document}
